% ============================================================================
% 01_modelo_dominio.tex — Cuerpo del entregable "Modelo de Dominio"
% Sin preámbulo: se incluye tanto desde 01_Modelo_Dominio.tex (standalone)
% como, eventualmente, desde un documento maestro.
% ============================================================================
\section{Modelo de Dominio}

\subsection{Propósito y alcance}

Esta sección presenta el \textbf{Modelo de Dominio} de LSA Salud siguiendo
la técnica descripta por Craig Larman en \textit{UML y Patrones} dentro del
Proceso Unificado. El modelo de dominio es un modelo conceptual: representa
los conceptos, atributos y asociaciones significativas del \emph{problema}
que LSA Salud resuelve (comunicación entre un paciente que usa Lengua de
Señas Argentina y personal de recepción o enfermería), no las clases de
software que lo implementan. Por esa razón, deliberadamente no aparecen aquí
elementos de diseño como \texttt{hooks} de React, conexiones WebSocket,
endpoints de FastAPI o archivos \texttt{.npy}: esos elementos pertenecen al
diseño del software y se documentan en otros entregables.

Las clases del diagrama no declaran métodos ni tipos de atributo —ambas
omisiones son intencionales en un modelo de dominio— y las asociaciones
llevan nombre y multiplicidad en ambos extremos, con una dirección de
lectura sugerida por el nombre del rol.

\subsection{Diagrama de clases conceptuales}

\diagramaApaisada{01_modelo_dominio}{0.95\textheight}{Modelo de Dominio de LSA Salud}{fig:modelo-dominio}

El diagrama se organiza en cuatro áreas conceptuales que se corresponden con
los tres actores de LSA Salud (ver Entregable de Casos de Uso): la
\textbf{comunicación} entre Paciente y Personal de Recepción o Enfermería
(parte izquierda e inferior del diagrama), el \textbf{reconocimiento de
señas} (parte central), la \textbf{gestión del modelo de reconocimiento} a
cargo del Desarrollador de IA / Equipo de Proyecto (parte derecha), y los
\textbf{conceptos diseñados y pendientes de implementación} que cubren los
requisitos oficiales RF-03, RF-04 y RF-06 (\textit{Consentimiento},
\textit{Auditoría} y \textit{Reporte de traducción}, ver
Sección~\ref{sec:dom-pendientes}).

\subsection{Diccionario de conceptos}

La Tabla~\ref{tab:diccionario-conceptos} describe cada concepto del modelo
de dominio, su rol en el problema y sus atributos conceptuales (sin tipo,
como corresponde a un modelo de dominio).

\begin{longtable}{p{3.6cm} p{5.4cm} p{5.4cm}}
    \toprule
    \textbf{Concepto} & \textbf{Descripción} & \textbf{Atributos} \\
    \midrule
    \endhead
    Paciente &
    Persona sorda o hipoacúsica que se comunica con el personal de salud
    mediante LSA frente a la cámara de la Estación del Paciente. &
    Sin datos persistentes: el sistema no almacena identidad ni historia
    clínica del paciente; se lo identifica únicamente por la sesión activa
    de su estación. \\

    Personal de Recepción o Enfermería &
    Personal de recepción o enfermería que lee la traducción del paciente
    y le responde por texto o por voz. &
    Sin datos persistentes (igual que Paciente). \\

    Desarrollador de IA / Equipo de Proyecto &
    Personal técnico del equipo que amplía el vocabulario de señas
    reconocibles y reentrena el modelo de reconocimiento. &
    Sin datos persistentes. \\

    Estación\-De\-Atención &
    Pantalla o dispositivo desde el que un actor humano interactúa con
    LSA Salud; existen dos por sesión de atención (una por paciente, una
    por personal de recepción o enfermería) y permanecen sincronizadas
    entre sí. &
    tipo (identifica si es la estación del paciente o la de recepción). \\

    Conversación &
    Intercambio de mensajes correspondiente a una sesión de atención entre
    un paciente y el personal de recepción o enfermería. &
    fechaInicio. \\

    Mensaje &
    Unidad de comunicación dentro de una conversación, emitida por el
    paciente (como resultado de una traducción) o por el personal de
    recepción o enfermería (escrito o dictado). &
    texto, rolEmisor, horaEnvío. \\

    Secuencia\-De\-Señas &
    Secuencia de palabras que el paciente va formando con señas sueltas,
    mientras la oración todavía está en curso (antes de traducirse a
    lenguaje natural). &
    estado (en curso / cerrada), marcaTemporalInicio. \\

    Clase\-De\-Seña &
    Categoría de seña que el sistema es capaz de reconocer (por ejemplo,
    ``dolor'', ``cabeza'', ``respirar'', ``yo'' o el estado de reposo
    ``idle''). &
    nombre. \\

    Palabra\-Reconocida &
    Ocurrencia concreta de una clase de seña, detectada dentro de una
    secuencia de señas, con su nivel de confianza. &
    confianza, posición (orden dentro de la secuencia), marcaTemporal. \\

    Traducción &
    Resultado de convertir una secuencia de señas cerrada en una oración
    en español natural. &
    textoOriginal (palabras sueltas), textoTraducido, fuente (plantilla,
    caché o modelo de lenguaje). \\

    Plantilla\-Médica &
    Frase médica predefinida, asociada a una combinación específica y
    conocida de clases de seña, que permite traducir sin recurrir a un
    modelo de lenguaje. &
    combinaciónDeSeñas, fraseResultante. \\

    Síntesis\-De\-Voz &
    Locución sonora generada a partir del texto de un mensaje, con acento
    argentino. &
    voz, formatoAudio. \\

    Modelo\-De\-Reconocimiento &
    Red neuronal entrenada capaz de clasificar una secuencia de
    movimientos en una clase de seña. &
    versión, cantidadDeClases, exactitudDeValidación. \\

    Muestra\-De\-Entrenamiento &
    Grabación de una seña realizada por el Desarrollador de IA / Equipo de
    Proyecto, utilizada para entrenar o reentrenar el modelo de
    reconocimiento. &
    cantidadDeFrames, fechaCaptura, modoDeCaptura (legado o basado en
    MediaPipe Tasks). \\
    \bottomrule
    \caption{Diccionario de conceptos del dominio de LSA Salud}
    \label{tab:diccionario-conceptos}
\end{longtable}

\subsection{Conceptos diseñados, pendientes de implementación}
\label{sec:dom-pendientes}

Los siguientes tres conceptos corresponden a los requisitos oficiales RF-03,
RF-04 y RF-06 de la matriz de trazabilidad (RF-03.1, RF-04.1 y RF-06.1 en
\texttt{IDS.md}), marcados como \textbf{diseñado --- pendiente de
implementación}: hoy no existe código que los materialice, pero se modelan
como parte del dominio para que su diseño quede completo desde esta primera
fase. La Tabla~\ref{tab:diccionario-pendientes} los describe con el mismo
nivel de detalle que el resto del diccionario.

\begin{longtable}{p{3.6cm} p{5.4cm} p{5.4cm}}
    \toprule
    \textbf{Concepto} & \textbf{Descripción} & \textbf{Atributos} \\
    \midrule
    \endhead
    Consentimiento\-De\-Cámara &
    Autorización explícita que el paciente otorga antes de que la Estación
    del Paciente active la cámara; condición obligatoria para que exista
    cualquier Secuencia\-De\-Señas (RF-03 oficial). &
    otorgado (booleano), marcaTemporal. \\

    Evento\-De\-Auditoría &
    Registro anónimo de una operación de sistema relevante (predicción,
    traducción, síntesis, respuesta, captura), sin datos identificatorios
    ni imágenes o video, generado automáticamente para auditoría operativa
    (RF-04 oficial). &
    tipoDeOperación, marcaTemporal, resultado (éxito o error). \\

    Reporte\-De\-Traducción &
    Señalamiento de que un mensaje de la conversación es una traducción
    incorrecta o ambigua, hecho por el Paciente o por el Personal de
    Recepción o Enfermería, para registrar el caso y mejorar el modelo
    (RF-06 oficial). &
    motivo, marcaTemporal. \\
    \bottomrule
    \caption{Diccionario de conceptos pendientes de implementación}
    \label{tab:diccionario-pendientes}
\end{longtable}

\subsection{Asociaciones principales}

La Tabla~\ref{tab:asociaciones-dominio} resume las asociaciones nombradas
del diagrama, con su multiplicidad y una lectura sugerida.

\begin{longtable}{p{3.4cm} p{3.6cm} p{3.4cm} p{3.6cm}}
    \toprule
    \textbf{Origen} & \textbf{Asociación} & \textbf{Destino} & \textbf{Lectura} \\
    \midrule
    \endhead
    Paciente & utiliza & Estación\-De\-Atención & Un paciente utiliza una estación de atención. \\
    Personal de Recepción o Enfermería & utiliza & Estación\-De\-Atención & El personal de recepción o enfermería utiliza una estación de atención. \\
    Estación\-De\-Atención & muestra & Conversación & Las dos estaciones de una sesión muestran la misma conversación. \\
    Paciente & otorga & Consentimiento\-De\-Cámara & Un paciente otorga el consentimiento antes de poder realizar señas (diseñado, pendiente). \\
    Paciente & realiza & Secuencia\-De\-Señas & Un paciente realiza cero o más secuencias de señas a lo largo de la sesión, condicionado a tener un Consentimiento\-De\-Cámara otorgado. \\
    Secuencia\-De\-Señas & compuesta por & Palabra\-Reconocida & Una secuencia de señas está compuesta por una o más palabras reconocidas. \\
    Palabra\-Reconocida & corresponde a & Clase\-De\-Seña & Cada palabra reconocida corresponde a una clase de seña del catálogo. \\
    Secuencia\-De\-Señas & da origen a & Traducción & Una secuencia de señas cerrada da origen, como máximo, a una traducción. \\
    Traducción & puede basarse en & Plantilla\-Médica & Una traducción puede basarse en una plantilla médica preexistente. \\
    Traducción & genera & Mensaje & Toda traducción genera exactamente un mensaje del paciente. \\
    Mensaje & tiene & Síntesis\-De\-Voz & Un mensaje puede tener una síntesis de voz asociada. \\
    Mensaje & pertenece a & Conversación & Todo mensaje pertenece a una conversación. \\
    Mensaje & puede tener & Reporte\-De\-Traducción & Un mensaje puede tener cero o un reporte de traducción incorrecta (diseñado, pendiente). \\
    Paciente & señala & Reporte\-De\-Traducción & Un paciente puede señalar cero o más reportes (diseñado, pendiente). \\
    Personal de Recepción o Enfermería & señala & Reporte\-De\-Traducción & El personal puede señalar cero o más reportes (diseñado, pendiente). \\
    Personal de Recepción o Enfermería & redacta & Mensaje & El personal de recepción o enfermería redacta cero o más mensajes. \\
    Desarrollador de IA / Equipo de Proyecto & captura & Muestra\-De\-Entrenamiento & Un desarrollador captura cero o más muestras de entrenamiento. \\
    Muestra\-De\-Entrenamiento & ejemplifica & Clase\-De\-Seña & Cada muestra ejemplifica una clase de seña. \\
    Desarrollador de IA / Equipo de Proyecto & entrena & Modelo\-De\-Reconocimiento & Un desarrollador entrena cero o más versiones del modelo. \\
    Modelo\-De\-Reconocimiento & reconoce & Clase\-De\-Seña & Un modelo reconoce una o más clases de seña. \\
    Modelo\-De\-Reconocimiento & clasifica & Palabra\-Reconocida & Un modelo clasifica cero o más palabras reconocidas. \\
    Estación\-De\-Atención & origina & Evento\-De\-Auditoría & Cada operación relevante de una estación origina cero o más eventos de auditoría anónimos (diseñado, pendiente). \\
    \bottomrule
    \caption{Asociaciones del Modelo de Dominio}
    \label{tab:asociaciones-dominio}
\end{longtable}
